\section{Experimental Design}
\subsection{Evaluation Objectives and Analysis Plan}
The evaluation focused primarily on task completion time under joint parameter-bundle configuration relative to the planned baselines. Success rate and Raw NASA-TLX were reported as secondary descriptive evidence of task success and subjective workload. Master-side trajectory length and pause count were included as exploratory process measures. Mode A served as the fixed-parameter baseline; mode B was a manual-selection workflow baseline, in which the operator identified the object and selected a strategy after timing began; and mode D was the visual-semantic display with fixed parameters. Mode E implemented impedance-only reconfiguration, retaining visual-semantic mapping and slave-side impedance updates but not the haptic-interface or gripper settings. The primary C--E comparison therefore evaluated the complete parameter bundle at the system level against impedance-only reconfiguration; it was not designed to separate the independent contributions of the added parameter groups or to test an interaction between them. All statistical comparisons are limited to repeated measurements in the present three-operator sample.

\subsection{Operators and Test Objects}
Three operators completed the main experiment (P01--P03; 23--24 years old; two male and one female; all right-handed). The Omega.7 was positioned to the operator's left and was operated with the left hand in every trial. All operators had received basic teleoperation training and completed a 10--15 min familiarization session immediately before formal data collection to practice master-device motion, gripper input, and the task procedure. All operators provided written informed consent.

The six objects were selected to represent different contact, grasping, and transfer requirements and were assigned to the fragility-priority, balanced, or stability-priority strategy (Table~\ref{tab:5}). These assignments were used solely to determine the parameter settings in this experiment and should not be interpreted as a general classification of material properties.

\begin{table*}[!t]
\centering
\caption{Test objects, assigned strategies, and principal task risks.}
\label{tab:5}
\small
\setlength{\tabcolsep}{6pt}
\renewcommand{\arraystretch}{1.16}
\begin{tabularx}{\textwidth}{@{}p{0.14\textwidth}p{0.18\textwidth}p{0.29\textwidth}Y@{}}
\toprule
Object & Assigned strategy & Physical profile & Principal task risk \\
\midrule
Apple & Fragility-priority & $\sim$200 g; smooth; $\varnothing$70--80 mm & Impact or slip; gentle contact required. \\
Banana & Fragility-priority & $\sim$120 g; smooth; $20\times180$ mm & Compression-induced deformation and slip. \\
Paper cup & Balanced & $\sim$5 g; paper; $\varnothing$75$\times$90 mm & Easily deformed; stable grasp required. \\
Bottle & Balanced & $\sim$30 g; smooth plastic; $\varnothing$65$\times$200 mm & Slip; balance between efficiency and stability. \\
Mouse & Stability-priority & $\sim$100 g; smooth plastic; $65\times120\times35$ mm & Rigid irregular surface; transfer slip. \\
Scissors & Stability-priority & $\sim$150 g; metal and plastic; $50\times170\times15$ mm & Rigid elongated geometry; high orientation demand. \\
\bottomrule
\end{tabularx}
\end{table*}
\subsection{Experimental Modes and Trial Structure}
The five modes and their comparison logic are defined in Table~\ref{tab:3}. Modes C, D, and E used the same visual-semantic display. The C--D contrast therefore compares joint parameter-bundle configuration with the visual-semantic display with fixed parameters, whereas C--E compares the full bundle with impedance-only reconfiguration while holding the designed interface constant. The D--A contrast assesses the complete visual-semantic display---class, mapped strategy, confidence, and bounding box---and does not isolate the class label. Mode B required manual object identification and button-based strategy selection after timing had started but before robot motion could begin. Its comparison with mode C therefore includes both automation of the selection workflow and differences in visual-interface presentation; the separate decision time was not logged and this contrast is not interpreted as an isolated controller effect.

The experiment comprised 27 matched task blocks. Each block was defined by a fixed operator, object, and repetition index and contained one trial in each of modes A--E, yielding $27\times5=135$ trials. Within each operator and strategy category, the two objects were assigned to the three repetitions in a fixed 2:1 ratio. This assignment was held constant across the five modes for a given operator and strategy, preserving the matched comparison, while the 2:1 allocation was reversed across operators. The resulting object counts yielded the approximate 5:4 balance within each strategy shown in Table~\ref{tab:6}.

\begin{table}[!t]
\centering
\caption{Allocation of matched task blocks and trials by strategy and object.}
\label{tab:6}
\footnotesize
\setlength{\tabcolsep}{6pt}
\renewcommand{\arraystretch}{1.16}
\begin{tabular*}{\columnwidth}{@{\extracolsep{\fill}}lcc@{}}
\toprule
Object & Blocks & Trials \\
\midrule
\multicolumn{3}{@{}l}{\textit{Fragility-priority}} \\
Apple & 4 & 20 \\
Banana & 5 & 25 \\
\addlinespace[2pt]
\multicolumn{3}{@{}l}{\textit{Balanced}} \\
Paper cup & 5 & 25 \\
Bottle & 4 & 20 \\
\addlinespace[2pt]
\multicolumn{3}{@{}l}{\textit{Stability-priority}} \\
Mouse & 5 & 25 \\
Scissors & 4 & 20 \\
\midrule
Total (six objects) & 27 & 135 \\
\bottomrule
\end{tabular*}
\par\smallskip
\footnotesize\textit{Note:} Each block contains one trial in each of the five modes.
\end{table}
The five mode orders were preassigned at the operator-by-strategy group level using nine different sequences, so no single fixed order was used. The schedule distributed modes across early and late positions but did not constitute a complete Latin-square design. Across the nine sequences, mode C occurred in position four in three groups, whereas mode E occurred in position five in three groups. The exact sequences are reported in Supplementary Table S2. With only three operators, order, object, and operator effects could not be fully separated; the C--E difference is therefore interpreted as a within-sample paired result rather than an order-free population effect.

\subsection{Task and Procedure}
Before each trial, the test object was manually placed within a predefined work region and at an approximately prescribed orientation. Small reset errors in position and orientation were therefore unavoidable, but the same placement protocol was used in all modes. Only translational master motion was mapped; the desired Panda end-effector orientation remained fixed throughout the trial.

Each trial comprised reset, approach, grasp, transfer, release, and termination (Fig.~\ref{fig:3}). In modes C--E, camera acquisition, visual inference, and teleoperation began concurrently at task initiation; visual output did not gate operator motion. In mode B, timing began before object identification and strategy selection, and robot motion was permitted only after the operator completed the button selection. No separate timestamp marked the end of this selection step.

A trial was classified as successful when grasping, transfer, and placement were completed without a drop, observable slip, or visible damage. Complete separation of the object from the gripper was classified as a drop; visible relative motion between object and gripper during grasp or transfer was classified as slip; and a clear indentation, crushing deformation, or other visible damage was classified as damage. A designated researcher applied these predefined criteria in real time. Trials were not systematically video-recorded, and the labels were not independently or blindly re-evaluated.

To maintain a common timing endpoint, failure did not stop the timer immediately. Regrasping after a drop was not permitted; the failure was recorded, and the operator completed the prescribed return-and-termination procedure until the robot and gripper reached the predefined terminal state. The same endpoint was therefore applied to successful and failed trials. The system recorded the master-side trajectory, gripper input, active control parameters, task duration, and outcome. It did not retain an event-level history of visual classes, confidences, strategy-trigger timestamps, or synchronized trigger--contact timestamps. Consequently, the retained active-state record cannot retrospectively establish whether the intended strategy was selected at a verified pre-contact time in each trial.

\subsection{Evaluation Measures}
Task completion time was the primary endpoint. Success rate was reported descriptively, and master-side trajectory length and pause count were used as exploratory process measures. Subjective workload was assessed with Raw NASA-TLX \cite{ref27}. Each of the six dimensions was rated from 0 to 100, and their unweighted arithmetic mean was used as the Raw NASA-TLX score.

Each operator completed one questionnaire immediately after the three trials in a given strategy $\times$ mode condition. These trials combined the two objects assigned to that strategy in the operator-specific 2:1 ratio. The workload dataset therefore contained $3$ operators $\times3$ strategies $\times5$ modes $=45$ questionnaires, or nine questionnaire units per mode. The resulting scores represent strategy-level workload under each mode and do not support object-specific workload comparisons or population-level workload inference.

Visual performance was characterized in a separate controlled image test using class accuracy, class-to-strategy mapping accuracy, confidence, and per-image processing time. A pause was defined as a continuous interval during which master-device speed remained below 0.005 m/s for at least 0.30 s. Pauses were detected by differentiating the raw master-side trajectories, which were sampled at approximately 200 Hz.

\subsection{Statistical Analysis}
Completion time across the five modes was first compared using a Friedman test, with the task block---defined by the same operator, object, and repetition index---as the pairing unit. When the omnibus test was significant, paired Wilcoxon signed-rank tests were performed with Holm--Bonferroni correction for multiple comparisons. For the primary C--E comparison, we report the paired mean difference, relative change, and a 95\% bootstrap confidence interval based on 10,000 resamples of the matched task blocks. We also examined effect direction for each operator and performed a leave-one-operator-out analysis to assess whether a single operator dominated the overall difference. The paired C--E success outcomes were tabulated, and an exact McNemar test was reported as a supplementary analysis because few discordant pairs were observed. Raw NASA-TLX scores were paired descriptively across nine operator $\times$ strategy units per mode; each score summarized three mixed-object trials, and no bootstrap interval was calculated. Trajectory length and pause count were treated as exploratory measures. Because the 27 task blocks were nested within only three operators, the inferential tests characterize repeated-measures differences within the present sample and do not support population-level inference across operators. Results are reported as both median [IQR] and mean $\pm$ SD. Figures~\ref{fig:4}--\ref{fig:9} were generated from frozen data using verified Python/Matplotlib scripts.
